\section{Datos del entorno y de la instalación}

La práctica se realizó en una estación de trabajo con Fedora Linux 44, una
distribución GNU/Linux de 64 bits. El equipo utilizó el núcleo
\texttt{7.1.13-200.fc44.x86\_64}. Antes de efectuar cambios se inspeccionó el
estado del sistema para distinguir entre componentes ausentes y componentes ya
instalados. Esta comprobación evitó reemplazar innecesariamente una
instalación funcional de Docker.

\begin{table}[H]
  \centering
  \rowcolors{2}{unamblue!6}{white}
  \begin{tabularx}{\textwidth}{>{\bfseries}lX}
    \toprule
    Componente & Versión registrada \\
    \midrule
    Sistema operativo & Fedora Linux 44 (Workstation Edition), x86\_64 \\
    Docker Engine y cliente & 29.7.2 \\
    containerd & 2.3.4 \\
    Docker Buildx & 0.36.1 \\
    PostgreSQL & 18.6, contenido en la imagen \texttt{postgres:18} \\
    DBeaver Community & 26.2.0 \\
    Controlador JDBC & PostgreSQL JDBC 42.7.13 \\
    \bottomrule
  \end{tabularx}
  \caption{Resumen del entorno utilizado durante la práctica.}
\end{table}

\evidencia{01_evidencia.png}{Identificación del sistema, paquetes instalados y estado del servicio Docker.}

El tiempo transcurrido entre la primera comprobación del entorno y la conexión
correcta desde DBeaver fue de aproximadamente cinco minutos. Este valor
corresponde al trabajo activo de habilitación, descarga, creación y prueba; no
incluye el tiempo posterior destinado a organizar capturas, redactar la
bitácora ni revisar el documento final.

\begin{nota}{Criterio de seguridad}
La contraseña de PostgreSQL se generó de manera aleatoria y se almacenó
temporalmente en un archivo con permisos restrictivos. Por esta razón, la
credencial no aparece en comandos, capturas ni fragmentos publicados.
\end{nota}
